\begin{frame}[plain]
  \sectiontag{Applied Forecasting Methods}
  \vspace{0.2em}
  {\fontsize{24}{28}\selectfont\bfseries Weighted Panel Forecasting on the Kaggle Hedge Fund Dataset\par}
  \vspace{0.2em}
  {\large\color{DeckMuted}Introduction and exploratory data analysis\par}
  \vspace{0.55em}

  \begin{columns}[T,onlytextwidth]
    \column{0.58\textwidth}
    \begin{softcard}
      {\bfseries What this PDF covers}\par
      \vspace{0.3em}
      This section introduces the forecasting task, explains the panel structure and evaluation logic, and uses EDA to justify what kinds of models are defensible later.
      \divider
      {\small
      \textbf{Scope.} Only the introduction and EDA story is included here. Classical and deep model results are intentionally left for the next section of the project.}
    \end{softcard}

    \column{0.38\textwidth}
    \begin{whitecard}
      \footnotesize
      {\bfseries Team}\par
      \TeamNames

      \vspace{0.35em}
      {\bfseries Roll Numbers}\par
      \RollNumbers

      \vspace{0.35em}
      {\bfseries Course}\par
      \CourseName

      \vspace{0.35em}
      {\bfseries Faculty}\par
      \FacultyName

      \vspace{0.35em}
      {\bfseries Institute}\par
      \InstituteName

      \vspace{0.35em}
      {\bfseries Date}\par
      \DeckDate
    \end{whitecard}
  \end{columns}
\end{frame}

\begin{frame}[t]{Roadmap}
  \sectiontag{Overview}
  \takeaway{The flow is simple: define the problem, define one series correctly, explain the split and metric, then use EDA to narrow the later modeling space.}
  \vspace{0.35em}

  \begin{center}
    \begin{tikzpicture}[x=1cm,y=1cm]
      \foreach \x/\n/\label in {
        0/1/Problem,
        2.45/2/Series Key,
        4.9/3/Data Scope,
        7.35/4/Split + Metric,
        9.8/5/EDA Findings,
        12.25/6/Implications
      }{
        \fill[DeckAccent] (\x,0) circle (0.24);
        \node[text=white,font=\bfseries\small] at (\x,0) {\n};
        \node[align=center,font=\bfseries\scriptsize,text=DeckInk] at (\x,-0.78) {\label};
      }
      \foreach \x/\nx in {0.38/2.05,2.83/4.5,5.28/6.95,7.73/9.4,10.18/11.85}{
        \draw[-{Latex[length=2.8mm]},line width=0.9pt,color=DeckAccent] (\x,0) -- (\nx,0);
      }
    \end{tikzpicture}
  \end{center}

  \vspace{0.45em}
  \begin{columns}[T,onlytextwidth]
    \column{0.48\textwidth}
    \begin{softcard}
      {\bfseries Focus of this PDF}\par
      \vspace{0.25em}
      \small
      \softbullet{Why the task is a weighted panel problem instead of one ordinary time series.}
      \softbullet{Why chronology, horizon, and series identity all matter at the same time.}
      \softbullet{Which EDA findings directly change later modeling decisions.}
    \end{softcard}

    \column{0.48\textwidth}
    \begin{warmcard}
      {\bfseries Not covered yet}\par
      \vspace{0.25em}
      \small
      \softbullet{No classical forecast comparison in this PDF.}
      \softbullet{No deep-learning result comparison in this PDF.}
      \softbullet{This closes exactly at the handoff from EDA into modeling.}
    \end{warmcard}
  \end{columns}
\end{frame}

\begin{frame}[t]{Problem Statement}
  \sectiontag{Overview}
  \takeaway{This is a weighted panel forecasting problem: many related series evolve over time, and their forecast errors do not contribute equally to the final score.}
  \vspace{0.35em}

  \begin{columns}[T,onlytextwidth]
    \column{0.48\textwidth}
    \begin{whitecard}
      {\bfseries Task definition}\par
      \vspace{0.25em}
      \small
      \begin{itemize}
        \item predict \inlinecode{y\_target} at future \inlinecode{ts\_index} values
        \item do this for many anonymized series, not just one
        \item each row carries a competition \inlinecode{weight}
        \item forecasts are evaluated on later timestamps, not shuffled rows
      \end{itemize}
    \end{whitecard}

    \column{0.48\textwidth}
    \begin{softcard}
      {\bfseries Why this is harder than one series}\par
      \vspace{0.25em}
      \small
      \begin{itemize}
        \item different series have different lengths, volatility, and scale
        \item the target is heavy-tailed, so rare moves matter
        \item weights are extremely concentrated, so some errors dominate
        \item horizons \inlinecode{1, 3, 10, 25} define different forecasting regimes
      \end{itemize}
    \end{softcard}
  \end{columns}
\end{frame}

\begin{frame}[t]{How One Series Is Defined}
  \sectiontag{Overview}
  \takeaway{One exact series is identified by \inlinecode{code + sub\_code + sub\_category + horizon}. \inlinecode{ts\_index} then orders the observations inside that series.}
  \vspace{0.22em}

  \begin{columns}[T,onlytextwidth]
    \column{0.57\textwidth}
    \begin{whitecard}
      {\bfseries Series-key schema}\par
      \vspace{0.22em}
      \begin{center}
        \begin{tikzpicture}[
          node distance=0.28cm,
          keybox/.style={draw=DeckBorder, rounded corners=3pt, fill=DeckSoft, minimum width=1.95cm, minimum height=0.78cm, font=\footnotesize\bfseries, align=center},
          flow/.style={-{Latex[length=2.4mm]}, line width=0.9pt, color=DeckAccent}
        ]
          \node[keybox] (code) {code};
          \node[keybox, right=of code] (subcode) {sub\_code};
          \node[keybox, right=of subcode] (subcat) {sub\_category};
          \node[keybox, right=of subcat] (horizon) {horizon};
          \node[draw=DeckAccent!25, rounded corners=3pt, fill=DeckWhite, minimum width=4.45cm, minimum height=0.82cm, below=0.6cm of subcode.east, anchor=north west, font=\small\bfseries, align=center] (series) {One unique series};
          \draw[flow] (code.south) -- ++(0,-0.35) -| (series.north west);
          \draw[flow] (subcode.south) -- ++(0,-0.25) -- (series.north);
          \draw[flow] (subcat.south) -- ++(0,-0.35) -| (series.north east);
          \draw[flow] (horizon.south) -- ++(0,-0.25) -| ([xshift=1.55cm]series.north east);

          \node[draw=DeckBorder, rounded corners=3pt, fill=DeckWhite, minimum width=5.5cm, minimum height=0.92cm, below=0.8cm of series, font=\scriptsize, align=center] (rows) {\textbf{Inside that series:} ordered by \inlinecode{ts\_index}\\each row stores \inlinecode{y\_target}, \inlinecode{weight}, and \inlinecode{feature\_*}};
          \draw[flow] (series.south) -- (rows.north);
        \end{tikzpicture}
      \end{center}
    \end{whitecard}

    \column{0.39\textwidth}
    \begin{softcard}
      {\bfseries What each field means}\par
      \vspace{0.16em}
      \scriptsize
      \textbf{Example key}\par
      \inlinecode{SJZP0OVU / OYJGNSQK / NQ58FVQM / H25}
      \vspace{0.2em}
      \begin{tabularx}{\linewidth}{L{0.35\linewidth}Y}
        \toprule
        \textbf{Field} & \textbf{Meaning} \\
        \midrule
        \inlinecode{code} & primary identifier \\
        \inlinecode{sub\_code} & secondary identifier \\
        \inlinecode{sub\_category} & categorical grouping \\
        \inlinecode{horizon} & forecast horizon \\
        \inlinecode{ts\_index} & time order within a series \\
        \inlinecode{y\_target} & continuous target value \\
        \inlinecode{weight} & row importance in scoring \\
        \bottomrule
      \end{tabularx}
    \end{softcard}
  \end{columns}
\end{frame}

\begin{frame}[t]{Dataset Scope And Length Rules}
  \sectiontag{Overview}
  \takeaway{The panel is large, heterogeneous, and uneven in series history, so later modeling has to be constrained by explicit length rules rather than applied blindly to every series.}
  \vspace{0.22em}

  \begin{columns}[T,onlytextwidth]
    \column{0.24\textwidth}\statcard{Train rows}{5.34M}
    \column{0.24\textwidth}\statcard{Test rows}{1.45M}
    \column{0.24\textwidth}\statcard{Features}{86}
    \column{0.24\textwidth}\statcard{Total series}{36,923}
  \end{columns}

  \vspace{0.22em}
  \begin{columns}[T,onlytextwidth]
    \column{0.48\textwidth}
    \begin{whitecard}
      {\bfseries Panel coverage}\par
      \vspace{0.18em}
      \small
      \begin{tabularx}{\linewidth}{L{0.47\linewidth}Y}
        \toprule
        \textbf{Measure} & \textbf{Value} \\
        \midrule
        Codes & 23 \\
        Train sub-codes & 180 \\
        Test sub-codes & 47 \\
        Sub-categories & 5 \\
        Horizons & 1, 3, 10, 25 \\
        Train timeline & \inlinecode{1..3601} \\
        Test timeline & \inlinecode{3602..4376} \\
        \bottomrule
      \end{tabularx}
    \end{whitecard}

    \column{0.48\textwidth}
    \begin{softcard}
      {\bfseries Length and split rules used later}\par
      \vspace{0.18em}
      \small
      \softbullet{\textbf{Length} = number of train timestamps available for one exact series.}
      \softbullet{\textbf{Minimum series length} = 120 points.}
      \softbullet{\textbf{Minimum train length} = 96 points.}
      \softbullet{\textbf{Minimum validation length} = 24 points.}
      \softbullet{\textbf{Lookback} = 24 recent target values.}
      \divider
      \footnotesize These thresholds stop very short series from creating misleading later comparisons.
    \end{softcard}
  \end{columns}
\end{frame}

\begin{frame}[t]{Timeline, Split, And Metric}
  \sectiontag{Overview}
  \takeaway{Chronology is preserved twice: once at the global train-versus-test boundary, and again inside each eligible series through an 80/20 chronological train-validation split.}
  \vspace{0.18em}

  \begin{whitecard}
    {\bfseries Timeline view of the forecasting problem}\par
    \vspace{0.18em}
    \begin{center}
      \begin{tikzpicture}[x=1cm,y=1cm]
        \fill[DeckAccent!20] (0,0) rectangle (8.55,0.78);
        \fill[DeckAccent!42] (8.55,0) rectangle (11.2,0.78);
        \draw[DeckBorder] (0,0) rectangle (11.2,0.78);
        \node[align=center,font=\bfseries\scriptsize] at (4.275,0.39) {Global train partition\\\inlinecode{1..3601}};
        \node[align=center,font=\bfseries\scriptsize,text=white] at (9.875,0.39) {Held-out test\\\inlinecode{3602..4376}};

        \fill[DeckSoft] (0,-0.88) rectangle (6.7,-0.2);
        \fill[DeckAccent!26] (6.7,-0.88) rectangle (8.35,-0.2);
        \draw[DeckBorder] (0,-0.88) rectangle (8.35,-0.2);
        \node[align=center,font=\bfseries\scriptsize] at (3.35,-0.54) {Per-series chronological train\\80\% of an eligible series};
        \node[align=center,font=\bfseries\scriptsize] at (7.53,-0.54) {Validation\\20\%};
      \end{tikzpicture}
    \end{center}
  \end{whitecard}

  \vspace{0.22em}
  \begin{columns}[T,onlytextwidth]
    \column{0.48\textwidth}
    \begin{softcard}
      {\bfseries Why the split matters}\par
      \vspace{0.16em}
      \small
      \softbullet{Random folds would leak future information into earlier states.}
      \softbullet{The held-out test is a later time suffix, so training must respect order.}
      \softbullet{The 80/20 per-series split gives later notebooks one unified validation rule.}
    \end{softcard}

    \column{0.48\textwidth}
    \begin{whitecard}
      {\bfseries Weighted skill score}\par
      \vspace{0.06em}
      {\small
      \[
        \mathrm{score} =
        \sqrt{
          1 -
          \mathrm{clip}_{0,1}
          \left(
            \frac{\sum_i w_i (y_i - \hat{y}_i)^2}{\sum_i w_i y_i^2}
          \right)
        }
      \]}
      \footnotesize
      A higher score is better. The weights \inlinecode{w\_i} mean that some rows contribute far more to the final ranking than others.
      \divider
      \begin{columns}[T,onlytextwidth]
        \column{0.31\textwidth}\smallstat{Top 1\%}{64.2\%}
        \column{0.31\textwidth}\smallstat{Top 5\%}{92.6\%}
        \column{0.31\textwidth}\smallstat{Top 10\%}{98.3\%}
      \end{columns}
    \end{whitecard}
  \end{columns}
\end{frame}
